%==============================================================================
% CAPÍTULO 17 — REPRODUTIBILIDADE E LICENÇAS
% PARTE III — Engenharia de Software Odontológico com SDD e TDD
%==============================================================================
\chapter{Reprodutibilidade e Licenças}
\label{cap:reprodutibilidade}
\nivelquatro

A Parte~III do livro começou com a disciplina da especificação (Capítulo~12)
e do teste (Capítulo~13), passou por projetos completos (Capítulo~14),
interpretabilidade (Capítulo~15) e auditoria de código alheio
(Capítulo~16). Este capítulo fecha o ciclo abordando a \textbf{publicação
responsável}: como garantir que \textbf{seu} código seja reproduzível,
corretamente licenciado e adequadamente documentado para a comunidade
científica e clínica.

\section{Reprodutibilidade Computacional em Odontologia}

\subsection{A crise de reprodutibilidade}

Em 2016, a revista \textit{Nature} publicou uma pesquisa histórica:
\textbf{mais de 70\% dos cientistas já tentaram — e falharam — em
reproduzir experimentos de outros pesquisadores}. Na IA odontológica,
o problema é agravado por três fatores:

\begin{enumerate}
  \item \textbf{Datasets privados}: a maioria dos datasets odontológicos
    contém dados de pacientes e não pode ser compartilhada publicamente.
  \item \textbf{Ambientes não versionados}: ``funciona na minha máquina''
    é a frase mais perigosa da ciência computacional.
  \item \textbf{Hiperparâmetros não documentados}: learning rate, batch
    size, augmentações — pequenas variações produzem grandes diferenças.
\end{enumerate}

\citacaoativa{baker2016reproducibility}{10.1038/533452a}{%
  ``A reprodutibilidade é um dos pilares do método científico. Sem ela,
  a ciência se torna anedótica, e a confiança pública na pesquisa se
  erosão.''}

\indice{reprodutibilidade computacional} \indice{crise de reprodutibilidade}

\subsection{Os 5 pilares da reprodutibilidade}

Para um projeto de IA odontológica ser considerado reproduzível, ele
deve atender a cinco pilares:

\begin{enumerate}
  \item \textbf{Código aberto e versionado}: repositório Git público com
    tags de versão correspondentes a cada resultado publicado.

  \item \textbf{Ambiente determinístico}: \texttt{requirements.txt} com
    versões exatas (\texttt{numpy==1.24.3}, não \texttt{numpy>=1.20}) ou
    Dockerfile com imagem congelada.

  \item \textbf{Dataset documentado}: se não puder ser compartilhado,
    descreva exatamente como foi obtido, processado e dividido. Forneça
    um \textbf{Data Card} (Seção~17.4).

  \item \textbf{Scripts de treinamento completos}: não apenas o modelo
    treinado, mas o script que o produziu — incluindo todas as
    augmentações, callbacks e seeds aleatórias.

  \item \textbf{Métricas na ``vida real''}: reporte métricas no conjunto
    de teste (nunca no de treino), com intervalos de confiança e
    \textit{k-fold cross-validation}.
\end{enumerate}

\destaque{A seed aleatória (\texttt{random\_state=42}) é o detalhe mais
negligenciado — e mais importante — para reprodutibilidade. Sem ela,
duas execuções idênticas produzem resultados diferentes.}

\section{Licenças: MIT, Apache, CC BY, GPL — O Que Pode Usar?}

A escolha da licença é uma decisão \textbf{legal}, não apenas técnica.
A Tabela~\ref{tab:licenses-comparison} compara as licenças mais relevantes
para projetos odontológicos.

\begin{table}[H]
  \centering
  \caption{Comparação de licenças para software e datasets odontológicos.}
  \label{tab:licenses-comparison}
  \begin{tabularx}{\textwidth}{lccccX}
    \toprule
    \textbf{Licença} & \textbf{Uso com.} & \textbf{Modif.} & \textbf{Patente} & \textbf{Copyleft} & \textbf{Recomendado para} \\
    \midrule
    MIT & Sim & Sim & Não & Não & Código flexível \\
    Apache 2.0 & Sim & Sim & Sim & Não & Código com proteção \\
    GPLv3 & Sim & Sim & Sim & Forte & Software livre \\
    CC BY 4.0 & Sim & Sim & N/A & Não & Datasets, artigos \\
    CC BY-SA 4.0 & Sim & Sim & N/A & Sim & Datasets abertos \\
    CC BY-NC 4.0 & Não & Sim & N/A & Não & Uso acadêmico \\
    CC0 & Sim & Sim & N/A & Não & Domínio público \\
    \bottomrule
  \end{tabularx}
\end{table}

\indice{licenças de software} \indice{MIT} \indice{Apache} \indice{GPL} \indice{Creative Commons}

\subsection{Recomendações práticas}

\begin{itemize}
  \item \textbf{Para código}: MIT se quiser máxima adoção; Apache 2.0 se
    quiser proteção de patentes; GPLv3 se quiser garantir que derivados
    permaneçam livres.

  \item \textbf{Para datasets}: CC BY 4.0 (atribuição obrigatória) ou
    CC0 (domínio público). \textbf{Evite} CC BY-NC se quiser que empresas
    usem seu dataset — o ``NC'' (não comercial) bloqueia uso por startups
    e indústria.

  \item \textbf{Para modelos treinados}: os pesos de uma rede neural são
    considerados ``obra derivada'' do dataset de treinamento. Se o dataset
    for CC BY-NC, o modelo também será NC.

  \item \textbf{Para documentação}: CC BY 4.0 é o padrão recomendado pela
    comunidade científica.
\end{itemize}

\notaexplicativa{Um erro comum é publicar código no GitHub \textbf{sem}
arquivo LICENSE. Na ausência de licença explícita, os termos padrão do
GitHub aplicam-se — você concede direito de \textit{view} e \textit{fork},
mas \textbf{não} de uso comercial, modificação ou redistribuição. Se
você quer que outros usem seu código, \textbf{escolha uma licença}.}

\subsection{Exercício: Escolhendo a Licença Correta}

\begin{enumerate}
  \item Você é um pesquisador que quer máxima citação e não se importa
    com uso comercial do seu código. $\rightarrow$ \textbf{MIT}
  \item Você é uma startup que quer proteger seu algoritmo de ser
    patentado por concorrentes. $\rightarrow$ \textbf{Apache 2.0}
  \item Você quer garantir que qualquer melhoria no seu código permaneça
    aberta. $\rightarrow$ \textbf{GPLv3}
  \item Você está publicando um dataset de radiografias anonimizadas e
    quer atribuição. $\rightarrow$ \textbf{CC BY 4.0}
  \item Você está publicando material educacional (este livro!).
    $\rightarrow$ \textbf{CC BY-NC-SA 4.0}
\end{enumerate}

\section{Model Card e Data Card}

\subsection{Model Card: a ``bula'' do seu modelo}

Proposto por \citeonline{mitchell2019model}, o \textit{Model Card} é um
documento estruturado que descreve:

\begin{itemize}
  \item \textbf{Detalhes do modelo}: nome, versão, arquitetura, framework.
  \item \textbf{Uso pretendido}: tarefa clínica, população-alvo, setting.
  \item \textbf{Uso fora do escopo}: situações em que \textbf{não} deve
    ser usado.
  \item \textbf{Dados de treinamento}: fonte, tamanho, demografia, divisão.
  \item \textbf{Métricas de desempenho}: acurácia, sensibilidade,
    especificidade, com intervalos de confiança e estratificação.
  \item \textbf{Limitações}: vieses conhecidos, populações não validadas.
  \item \textbf{Licença e citação}: como citar e sob qual licença usar.
\end{itemize}

\subsection{Data Card: a ``certidão de nascimento'' do dataset}

Proposto por \citeonline{gebru2021datasheets}, o \textit{Data Card}
(ou \textit{Datasheet}) documenta:

\begin{itemize}
  \item \textbf{Motivação}: por que o dataset foi criado?
  \item \textbf{Composição}: quantas amostras, classes, distribuição.
  \item \textbf{Coleta}: período, instituição, equipamento utilizado.
  \item \textbf{Pré-processamento}: todas as transformações aplicadas.
  \item \textbf{Anonimização}: como os dados dos pacientes foram protegidos.
  \item \textbf{Aprovação ética}: número CAAE ou equivalente.
  \item \textbf{Uso recomendado e restrições}.
\end{itemize}

\begin{pratica}{Gerando Model Card e Data Card para publicação}
  Objetivo: Criar documentação completa para o classificador de cáries
  do Projeto 1 (Capítulo~14) usando templates estruturados.
\end{pratica}

\codigo{codigos/cap17/publish_model_hf.py}

\teste{O Model Card gerado contém todos os campos obrigatórios segundo
Mitchell et al. (2019): detalhes do modelo, uso pretendido, dados de
treinamento, métricas, limitações e licença. O Data Card inclui
motivação, composição, coleta, anonimização e aprovação ética.}

\section{Como Publicar Modelo no Hugging Face}

O Hugging Face Model Hub é a plataforma mais utilizada para compartilhar
modelos de IA. Publicar seu modelo odontológico lá oferece:

\begin{itemize}
  \item \textbf{DOI automático}: cada modelo recebe um DOI via Zenodo.
  \item \textbf{Versionamento}: cada commit gera uma versão rastreável.
  \item \textbf{API de inferência gratuita}: outros podem testar seu
    modelo sem instalar nada.
  \item \textbf{Model Card integrado}: documentação fica junto com o modelo.
  \item \textbf{Descoberta}: comunidade de 500k+ desenvolvedores.
\end{itemize}

\subsection{Passo a passo}

\begin{enumerate}
  \item \textbf{Crie uma conta} em \url{https://huggingface.co}.
  \item \textbf{Instale a biblioteca}:
\begin{lstlisting}[language=bash]
!pip install huggingface_hub
\end{lstlisting}

  \item \textbf{Autentique-se}:
\begin{lstlisting}[language=Python]
from huggingface_hub import notebook_login
notebook_login()  # Cole seu token de acesso
\end{lstlisting}

  \item \textbf{Salve e envie o modelo}:
\begin{lstlisting}[language=Python]
from huggingface_hub import HfApi, create_repo

# Cria repositório
repo_name = "edsonlaranjeiras/CariesICDAS-EfficientNetB3"
create_repo(repo_name, exist_ok=True)

# Envia modelo + Model Card
api = HfApi()
api.upload_file(
    path_or_fileobj="model.h5",
    path_in_repo="model.h5",
    repo_id=repo_name,
)
api.upload_file(
    path_or_fileobj="model_card.md",
    path_in_repo="README.md",
    repo_id=repo_name,
)
\end{lstlisting}

  \item \textbf{Adicione metadados} no arquivo \texttt{README.md} (YAML
    frontmatter):
\begin{lstlisting}[language=yaml]
---
license: mit
language: pt
tags:
  - odontologia
  - caries
  - classificacao
  - radiografia
  - bitewing
datasets:
  - edsonlaranjeiras/ICDAS-Brasil-2024
metrics:
  - accuracy
  - sensitivity
  - specificity
  - auc
pipeline_tag: image-classification
---
\end{lstlisting}

  \item \textbf{Ative o DOI via Zenodo}: nas configurações do repositório
    no Hugging Face, ative a integração com Zenodo.
\end{enumerate}

\teste{Modelo publicado no Hugging Face Hub com DOI automático,
Model Card integrado, tag \texttt{odontologia} e pipeline de
classificação de imagens configurado.}

\section{Como Citar GitHub + DOI + Dataset em ABNT}

A citação correta de software, datasets e modelos é uma exigência
crescente em periódicos odontológicos. A ABNT NBR 6023:2018 prevê
formatos específicos.

\subsection{Citação de software (ABNT NBR 6023)}

\begin{quote}
  \textbf{Formato:} SOBRENOME, Nome. \textbf{Título do software}
  [software]. Versão. Local: Distribuidor, Ano. Disponível em: URL.
  Acesso em: data.
\end{quote}

\begin{mdframed}[linecolor=corSecundaria, linewidth=1pt, backgroundcolor=corSecundaria!5, roundcorner=4pt]
  \textbf{Exemplo de citação ABNT:} \\
  LARANJEIRAS, Edson. \textbf{CariesICDAS-EfficientNetB3} [software].
  v1.0.0. Hugging Face Model Hub, 2026. Disponível em:
  \url{https://huggingface.co/edsonlaranjeiras/CariesICDAS-EfficientNetB3}.
  Acesso em: 09 jul. 2026.
\end{mdframed}

\subsection{Citação de dataset}

\begin{quote}
  \textbf{Formato:} SOBRENOME, Nome. \textbf{Título do dataset} [dataset].
  Versão. Plataforma, Ano. DOI: xx.xxxx/xxxxx.
\end{quote}

\subsection{Citação de repositório GitHub com DOI}

Quando o repositório tem DOI (via Zenodo), cite o DOI, não o link do
GitHub (links podem mudar, DOIs são permanentes):

\begin{mdframed}[linecolor=corSecundaria, linewidth=1pt, backgroundcolor=corSecundaria!5, roundcorner=4pt]
  \textbf{Exemplo de citação com DOI:} \\
  LARANJEIRAS, Edson. \textbf{OdontoIA: classificador de cáries ICDAS}
  [código-fonte]. v1.0.0. Zenodo, 2026. DOI: 10.5281/zenodo.1234567.
\end{mdframed}

\destaque{Sempre cite o \textbf{DOI} (via Zenodo), não o link do GitHub.
O GitHub pode alterar URLs, excluir repositórios ou tornar-se inacessível.
O DOI é permanente e rastreável.}

\section{Ética de Uso de Código Alheio}

\subsection{Atribuição é obrigatória — e é lei}

Usar código open source sem atribuição é \textbf{violação de licença}
— o equivalente jurídico a plágio. A licença MIT, a mais permissiva,
exige apenas que o aviso de copyright original seja mantido. Ignorar
esse requisito expõe você a riscos legais.

\subsection{Boa prática: dependências explícitas}

\begin{itemize}
  \item Liste \textbf{todas} as dependências no \texttt{requirements.txt}
    com versões exatas.
  \item No artigo, inclua uma seção ``Software Utilizado'' com citações
    ABNT de cada dependência relevante.
  \item Se você modificou código de terceiros, documente as modificações
    e mantenha os avisos de copyright originais.
\end{itemize}

\subsection{Checklist ético antes de publicar}

\begin{enumerate}
  \item \textbf{Licença}: todo arquivo de código tem uma licença compatível?
  \item \textbf{Atribuição}: código de terceiros está corretamente atribuído?
  \item \textbf{Dados de pacientes}: foram anonimizados? O CEP aprovou?
  \item \textbf{Viés}: o modelo foi testado em subgrupos demográficos?
  \item \textbf{Limitações}: estão claramente declaradas no Model Card?
  \item \textbf{Uso indevido}: você declarou usos fora do escopo?
\end{enumerate}

\section{Síntese do Capítulo e da Parte III}

Este capítulo concluiu a Parte~III do livro, cobrindo os aspectos legais
e éticos da publicação de software odontológico com IA:

\begin{itemize}
  \item \textbf{Reprodutibilidade} em 5 pilares: código aberto, ambiente
    determinístico, dataset documentado, scripts completos, métricas
    honestas.

  \item \textbf{Licenças}: MIT para código flexível, Apache 2.0 com
    proteção de patentes, CC BY 4.0 para datasets. \textbf{Nunca publique
    sem licença.}

  \item \textbf{Model Card e Data Card}: documentação estruturada que
    acompanha modelo e dataset, permitindo uso informado e responsável.

  \item \textbf{Hugging Face}: plataforma para publicação com DOI
    automático, versionamento e API de inferência gratuita.

  \item \textbf{Citação ABNT}: formatos específicos para software,
    dataset e repositório, sempre priorizando o DOI.

  \item \textbf{Ética}: atribuição é obrigatória; dados de pacientes
    exigem anonimização e aprovação ética; limitações e vieses devem
    ser declarados.
\end{itemize}

\subsection*{Balanço da Parte III}

A Parte~III percorreu o ciclo completo da engenharia de software
odontológico com qualidade profissional:

\begin{center}
  \begin{tikzpicture}[
    box/.style={rectangle, draw=corPrimaria, fill=corPrimaria!5, text width=3.5cm, align=center, align=center, minimum height=1.8cm, rounded corners=4pt, font=\footnotesize},
    arrow/.style={->, >=stealth, thick, corDestaque},
  ]
    \node[align=center, box] (spec) {\textbf{Cap 12 — SDD}\\\footnotesize Especificar\\\footnotesize antes de codar};
    \node[align=center, box, right=0.5cm of spec] (tdd) {\textbf{Cap 13 — TDD}\\\footnotesize Testar antes\\\footnotesize do modelo};
    \node[align=center, box, right=0.5cm of tdd] (proj) {\textbf{Cap 14 — Projetos}\\\footnotesize 4 projetos\\\footnotesize SDD+TDD};
    \node[align=center, box, below=1cm of tdd] (xai) {\textbf{Cap 15 — XAI}\\\footnotesize Interpretar\\\footnotesize e explicar};
    \node[align=center, box, left=0.5cm of xai] (audit) {\textbf{Cap 16 — Auditoria}\\\footnotesize Auditar código\\\footnotesize de terceiros};
    \node[align=center, box, right=0.5cm of xai] (reprod) {\textbf{Cap 17 — Publicação}\\\footnotesize Licenciar, documentar\\\footnotesize e publicar};

    \draw[arrow] (spec) -- (tdd);
    \draw[arrow] (tdd) -- (proj);
    \draw[arrow] (proj) -- (xai);
    \draw[arrow] (xai) -- (audit);
    \draw[arrow] (xai) -- (reprod);
    \draw[arrow, bend right=30] (reprod) to (spec);
  \end{tikzpicture}
\end{center}

A Parte~IV do livro aplicará todo esse arcabouço metodológico às
especialidades odontológicas: periodontia, implantodontia, ortodontia
e gêmeos digitais — sempre com SDD, TDD e XAI como alicerces de
qualidade.

\section{Ambientes Reproduzíveis com Docker e Conda}

\subsection{Docker: ``Funciona na minha máquina'' resolvido}

O Docker encapsula todo o ambiente de execução — sistema operacional,
bibliotecas, dependências, variáveis de ambiente — em um contêiner
imutável. Para projetos odontológicos, isso significa que qualquer
pesquisador, em qualquer sistema operacional, pode reproduzir seus
resultados com um único comando.

\begin{lstlisting}[language=docker, caption=Dockerfile para ambiente odontológico reproduzível]
FROM tensorflow/tensorflow:2.15.0-gpu

# Dependências do sistema
RUN apt-get update && apt-get install -y \
    libgl1-mesa-glx \
    libgomp1 \
    && rm -rf /var/lib/apt/lists/*

# Dependências Python com versões exatas
COPY requirements.txt /tmp/
RUN pip install --no-cache-dir -r /tmp/requirements.txt

# Código do projeto
COPY . /app
WORKDIR /app

# Comando padrão: treinar e avaliar
CMD ["python", "train.py", "--epochs", "100", "--batch-size", "32"]
\end{lstlisting}

Para construir e executar:

\begin{lstlisting}[language=bash]
docker build -t odontoia-caries:1.0.0 .
docker run --gpus all -v $(pwd)/data:/app/data odontoia-caries:1.0.0
\end{lstlisting}

\subsection{Conda: Reprodutibilidade para Ciência de Dados}

Para quem prefere o ecossistema Conda (comum em bioinformática e
imagens médicas), o \texttt{environment.yml} oferece controle preciso:

\begin{lstlisting}[language=yaml, caption=environment.yml com versões exatas]
name: odontoia
channels:
  - conda-forge
  - defaults
dependencies:
  - python=3.10.13
  - numpy=1.24.3
  - opencv=4.8.1
  - matplotlib=3.7.2
  - scikit-learn=1.3.0
  - pip
  - pip:
    - tensorflow==2.15.0
    - torch==2.1.0
    - pydicom==2.4.3
    - odontoia==1.0.0
\end{lstlisting}

\subsection{Pinning de Versões: O Detalhe que Faz a Diferença}

A diferença entre \texttt{numpy>=1.20} e \texttt{numpy==1.24.3} é a
diferença entre ``provavelmente funciona'' e ``certamente funciona''.
Sempre especifique versões \textbf{exatas}:

\begin{itemize}
  \item \textbf{Python}: \texttt{python==3.10.13}
  \item \textbf{NumPy}: \texttt{numpy==1.24.3}
  \item \textbf{TensorFlow}: \texttt{tensorflow==2.15.0}
  \item \textbf{PyTorch}: \texttt{torch==2.1.0}
  \item \textbf{CUDA/cuDNN}: compatível com a versão do framework
\end{itemize}

\section{Versionamento de Dados com DVC}

O Git é excelente para código, mas péssimo para datasets grandes (50 GB
de radiografias DICOM). O \textbf{DVC} (Data Version Control) resolve
isso versionando datasets em armazenamento externo (S3, GCS, SSH) e
mantendo apenas ponteiros leves no Git.

\begin{lstlisting}[language=bash, caption=Workflow DVC para dataset odontológico]
# Inicializa DVC no projeto
dvc init

# Adiciona dataset ao controle de versão
dvc add data/radiografias/
git add data/radiografias.dvc data/.gitignore
git commit -m "Dataset: 3.400 radiografias bitewing (ICDAS-Brasil-2024)"

# Configura armazenamento remoto (ex: Google Drive)
dvc remote add -d storage gdrive://1a2b3c4d5e6f
dvc push

# Outro pesquisador recupera o dataset exato:
git clone https://github.com/usuario/repo.git
dvc pull  # Baixa os dados do armazenamento remoto
\end{lstlisting}

\section{Publicação em Zenodo e Atribuição de DOI}

O \textbf{Zenodo} (\url{https://zenodo.org}) é um repositório europeu
de acesso aberto que atribui DOIs gratuitos a software, datasets e
artigos. Ele integra-se nativamente com o GitHub: cada release gera
automaticamente um registro no Zenodo com DOI.

\subsection{Passo a Passo GitHub + Zenodo}

\begin{enumerate}
  \item Acesse \url{https://zenodo.org} e faça login com sua conta GitHub.
  \item Autorize o Zenodo a acessar seus repositórios.
  \item No Zenodo, ative o \textit{webhook} para o repositório desejado.
  \item No GitHub, crie uma \textbf{release} (ex.: \texttt{v1.0.0}).
  \item Automaticamente, o Zenodo criará um registro com:
    \begin{itemize}
      \item DOI único (ex.: \texttt{10.5281/zenodo.1234567})
      \item Metadados extraídos do repositório (autores, descrição, licença)
      \item Badge de DOI para adicionar ao README
    \end{itemize}
\end{enumerate}

\destaque{Sempre que publicar um artigo que usa seu código, \textbf{faça
uma release no GitHub antes de submeter o artigo}. O DOI estará disponível
para incluir na seção ``Disponibilidade de Código'' do manuscrito.}

\section{Checklist Final de Publicação}

Antes de publicar qualquer projeto odontológico de IA, verifique esta
lista exaustiva:

\subsection{ checklist — Código}

\begin{itemize}
  \item[\checkmark] Repositório público no GitHub/GitLab
  \item[\checkmark] Arquivo LICENSE compatível com o uso pretendido
  \item[\checkmark] README.md com: descrição, instalação, uso, exemplos, citação
  \item[\checkmark] requirements.txt ou environment.yml com versões exatas
  \item[\checkmark] Dockerfile para ambiente reproduzível
  \item[\checkmark] Scripts de treinamento completos (não apenas modelo salvo)
  \item[\checkmark] Seeds aleatórias fixadas (\texttt{random\_state=42})
  \item[\checkmark] Testes automatizados com pytest (cobertura $\geq$ 80\%)
  \item[\checkmark] CI/CD via GitHub Actions
  \item[\checkmark] Release no GitHub com tag semântica (v1.0.0)
  \item[\checkmark] DOI via Zenodo
\end{itemize}

\subsection{ checklist — Modelo}

\begin{itemize}
  \item[\checkmark] Model Card completo (nome, tarefa, arquitetura, dataset, métricas, limitações)
  \item[\checkmark] Pesos do modelo publicados (Hugging Face, Zenodo ou GitHub Releases)
  \item[\checkmark] Script de inferência independente (não requer re-treinamento)
  \item[\checkmark] Formatos de exportação: H5 ou SavedModel (TensorFlow), .pt (PyTorch), .onnx (intercâmbio)
  \item[\checkmark] Testes de inferência com dados sintéticos (para verificação rápida)
\end{itemize}

\subsection{ checklist — Dataset}

\begin{itemize}
  \item[\checkmark] Data Card completo (motivação, composição, coleta, anonimização)
  \item[\checkmark] Aprovação ética documentada (número CAAE ou equivalente)
  \item[\checkmark] Se público: DOI via Zenodo, Kaggle ou similar
  \item[\checkmark] Se privado: instruções claras para solicitação de acesso
  \item[\checkmark] Descrição do pré-processamento aplicado
  \item[\checkmark] Divisão treino/validação/teste documentada e estratificada
  \item[\checkmark] Estatísticas descritivas (média, desvio, balanceamento)
\end{itemize}

\section{Ética de Uso de Código Alheio — Casos Reais}

\subsection{Caso 1: O Perigo do Copy-Paste sem Licença}

Em 2023, um grupo de pesquisa brasileiro publicou um artigo descrevendo
um ``novo'' classificador de lesões bucais. Uma auditoria revelou que
85\% do código havia sido copiado de um repositório GPLv3 — mas o artigo
não mencionava a fonte, e o código derivado foi publicado sob licença
proprietária. Resultado: \textbf{retratação do artigo} e danos à
reputação do grupo.

\subsection{Caso 2: Dataset ``Anonimizado'' que Não Era}

Um dataset de radiografias panorâmicas publicado no Kaggle como ``CC0''
(domínio público) foi retirado do ar após pesquisadores descobrirem que
os metadados DICOM continham nomes de pacientes, datas de nascimento e
números de prontuário. A ``anonimização'' havia removido apenas os tags
mais óbvios, mas tags privados (\textit{private tags}) continham dados
identificáveis.

\notaexplicativa{Para anonimizar DICOM corretamente, use ferramentas como
\texttt{pydicom} com \texttt{dataset.remove\_private\_tags()} E verificação
manual. Automação sem verificação humana é insuficiente para conformidade
com a LGPD.}

\subsection{Lições Aprendidas}

\begin{enumerate}
  \item \textbf{Atribuição é lei, não cortesia.} Copiar código sem
    atribuição viola a licença e constitui plágio acadêmico.

  \item \textbf{Anonimização é processo, não evento.} Requer ferramentas
    adequadas, verificação humana e auditoria por pares.

  \item \textbf{Licença é decisão de projeto.} Defina a licença antes da
    primeira linha de código, não como \textit{afterthought}.
\end{enumerate}

\section{Síntese do Capítulo e da Parte III}

Este capítulo concluiu a Parte~III do livro, cobrindo os aspectos legais
e éticos da publicação de software odontológico com IA:

\begin{itemize}
  \item \textbf{Reprodutibilidade} em 5 pilares: código aberto, ambiente
    determinístico (Docker/Conda), dataset documentado (DVC), scripts
    completos, métricas honestas.

  \item \textbf{Licenças}: MIT para código flexível, Apache 2.0 com
    proteção de patentes, CC BY 4.0 para datasets. \textbf{Nunca publique
    sem licença.} CC BY-NC bloqueia uso comercial — escolha com cuidado.

  \item \textbf{Model Card e Data Card}: documentação estruturada que
    acompanha modelo e dataset, permitindo uso informado e responsável.

  \item \textbf{Hugging Face + Zenodo}: plataforma de publicação com
    DOI automático, versionamento e API de inferência gratuita.

  \item \textbf{Citação ABNT}: formatos específicos para software,
    dataset e repositório, sempre priorizando o DOI sobre URLs voláteis.

  \item \textbf{Docker e DVC}: ferramentas que garantem reprodutibilidade
    computacional independente de sistema operacional ou hardware.

  \item \textbf{Checklist de publicação}: 30+ itens verificáveis que
    garantem que seu projeto está pronto para a comunidade científica.

  \item \textbf{Casos reais de falhas éticas}: lições aprendidas com
    retratações, vazamento de dados e violações de licença.
\end{itemize}

\subsection*{Balanço da Parte III}

A Parte~III percorreu o ciclo completo da engenharia de software
odontológico com qualidade profissional — da especificação (Capítulo~12)
à publicação responsável (este capítulo). O diagrama abaixo sintetiza a
jornada:

\begin{center}
  \begin{tikzpicture}[
    box/.style={rectangle, draw=corPrimaria, fill=corPrimaria!5, text width=3.5cm, align=center, align=center, minimum height=1.8cm, rounded corners=4pt, font=\footnotesize},
    arrow/.style={->, >=stealth, thick, corDestaque},
  ]
    \node[align=center, box] (spec) {\textbf{Cap 12 — SDD}\\\footnotesize Especificar\\\footnotesize antes de codar};
    \node[align=center, box, right=0.5cm of spec] (tdd) {\textbf{Cap 13 — TDD}\\\footnotesize Testar antes\\\footnotesize do modelo};
    \node[align=center, box, right=0.5cm of tdd] (proj) {\textbf{Cap 14 — Projetos}\\\footnotesize 4 projetos\\\footnotesize SDD+TDD};
    \node[align=center, box, below=1cm of tdd] (xai) {\textbf{Cap 15 — XAI}\\\footnotesize Interpretar\\\footnotesize e explicar};
    \node[align=center, box, left=0.5cm of xai] (audit) {\textbf{Cap 16 — Auditoria}\\\footnotesize Auditar código\\\footnotesize de terceiros};
    \node[align=center, box, right=0.5cm of xai] (reprod) {\textbf{Cap 17 — Publicação}\\\footnotesize Licenciar, documentar\\\footnotesize e publicar};

    \draw[arrow] (spec) -- (tdd);
    \draw[arrow] (tdd) -- (proj);
    \draw[arrow] (proj) -- (xai);
    \draw[arrow] (xai) -- (audit);
    \draw[arrow] (xai) -- (reprod);
    \draw[arrow, bend right=30] (reprod) to (spec);
  \end{tikzpicture}
\end{center}

A Parte~IV do livro aplicará todo esse arcabouço metodológico às
especialidades odontológicas: periodontia, implantodontia, ortodontia
e gêmeos digitais — sempre com SDD, TDD e XAI como alicerces de
qualidade.
